\section{Negative energy, Morse index, and the momentum response matrix}
\label{sec:r10-inertia}

This section distinguishes two statements.  The negative-index obstruction
below holds for every exact profile and requires no nondegeneracy.  Statements
about derivatives with respect to velocity apply only when a differentiable
family of profiles has been supplied.  Such a family is not presumed to exist
at every profile.

Let $a$ be a real-even positive elliptic homogeneous symbol of degree $m>1$,
let $X=H^{m/2}(\mathbb R^N;\mathbb C)$, and set
\[
 \alpha=\frac{2m}{N},\qquad
 E(U)=\frac12\int a(\xi)|\widehat U(\xi)|^2\,d\xi
       -\frac1{\alpha+2}\int|U|^{\alpha+2},\qquad
 P(U)=\operatorname{Im}\int\overline U\nabla U.
\]
Suppose that a nonzero $U\in X$ solves
\begin{equation}\label{eq:r10-profile}
 a(D)U+\omega U+i v\cdot\nabla U-|U|^\alpha U=0.
\end{equation}
The mass-critical Pohozaev identity is $v\cdot P(U)=2mE(U)$.
All pairings below are form pairings on $X$ and $X^*$; in particular
the derivative pairings are defined since $m>1$.
Our complex pairing is antilinear in its first argument, and the real
Hilbert pairing is its real part.

Define the real scalar operator and the complex-linear operator
\[
 A_U=a(D)+\omega-|U|^\alpha,\qquad
 K_U=A_U+i v\cdot\nabla.
\]
Both are selfadjoint as forms on the complex space, and $A_U$ commutes
with complex conjugation.  The Hessian of
$S_{\omega,v}=E+\omega M/2-v\cdot P/2$ is a \emph{real-linear} operator
$\mathcal L_U$, whose quadratic form is
\begin{equation}\label{eq:r10-fullhessian}
 q_{\mathcal L_U}(\eta)
 =\operatorname{Re}\langle\eta,K_U\eta\rangle
 -\alpha\int |U|^{\alpha-2}
       \bigl(\operatorname{Re}(\overline U\eta)\bigr)^2.
\end{equation}
At $U=0$ the last integrand is defined to be zero.  It is bounded by
$|U|^\alpha|\eta|^2$ and defines a continuous form even when $\alpha<2$.
The notation $n_-(\mathcal L)$ denotes the maximal dimension of a real
subspace on which its quadratic form is strictly negative.  It does not
denote the number of exponentially growing modes of the evolution.

\subsection{An unconditional negative-index obstruction}

\begin{theorem}[Two scalar directions and three action directions]
\label{thm:r10-three-directions}
If $E(U)<0$, then
\begin{equation}\label{eq:r10-index-bound}
 n_-(A_U|_{X_{\mathbb R}})\ge2,
 \qquad n_-(\mathcal L_U)\ge3.
\end{equation}
More generally these conclusions require only $v\cdot P(U)<0$ and
the profile equation; homogeneity and criticality are used solely to
infer that sign from the energy.

Consequently, a profile whose full action Hessian has negative index
at most two cannot have negative energy.  These statements are uniform
in frequency and speed and do not require proximity to a ground state.
\end{theorem}
\begin{proof}
Write $U=h+ib$ with real $h,b$.  The profile equation gives
\[
 A_Uh=v\cdot\nabla b,\qquad A_Ub=-v\cdot\nabla h.
\]
Antisymmetry of the derivative and
$P(U)=-2\langle b,\nabla h\rangle$ imply
\begin{equation}\label{eq:r10-scalar-plane}
 \langle h,A_Uh\rangle=\langle b,A_Ub\rangle
       =\frac12v\cdot P(U),\qquad
 \langle h,A_Ub\rangle=0.
\end{equation}
The functions $h,b$ are linearly independent: dependence would make
$U$ a constant phase times a real function and would force $P(U)=0$.
Thus their real span is a two-dimensional negative subspace.

For the full Hessian it is useful to retain complex notation.
Since $K_UU=0$, conjugation of the profile equation gives
\[
 K_U\overline U=2i v\cdot\nabla\overline U,
 \qquad
 \langle\overline U,K_U\overline U\rangle=2v\cdot P(U).
\]
The latter pairing is real.  Selfadjointness and $K_UU=0$ show
that all $K_U$ form pairings between $U$ and $\overline U$ vanish.
Moreover, $U$ and $\overline U$ are linearly independent over
$\mathbb C$, for the same constant-phase reason as above.

For $t\in\mathbb R$ and $z\in\mathbb C$ put
$\eta=tU+z\overline U$.  Equation \eqref{eq:r10-fullhessian} becomes
\begin{equation}\label{eq:r10-three-plane}
 q_{\mathcal L_U}(tU+z\overline U)
 =2v\cdot P(U)|z|^2
 -\alpha\int |U|^{\alpha-2}
     \left(t|U|^2+\operatorname{Re}(z\overline{U}^{2})\right)^2.
\end{equation}
If $z\ne0$ the first term is strictly negative.  If $z=0$ and
$t\ne0$, the second term equals
$-\alpha t^2\|U\|_{\alpha+2}^{\alpha+2}<0$.
The real span of $U,\overline U,i\overline U$ is therefore a
three-dimensional negative subspace.  On the full four-dimensional
real space $\operatorname{span}_{\mathbb C}\{U,\overline U\}$,
the restricted form has precisely three negative directions and the
one-dimensional kernel spanned by $iU$.
\end{proof}

\begin{remark}[What the obstruction does and does not say]
The proof is an explicit finite-dimensional compression of the
Hessian, not a computation of its full spectrum.  If the essential
spectrum reaches negative values, the variational index may be
infinite; three discrete negative eigenvalues are not being asserted.
In particular it
does not by itself prove spectral or nonlinear instability.  The same
argument applies to a real dispersion operator and a gauge-invariant
nonlinearity $f(|u|^2)u$ whenever $f'\ge0$ on the range of $|U|^2$
and $\int f'(|U|^2)|U|^4>0$: the last term in
\eqref{eq:r10-fullhessian} is then replaced by
$-2\int f'(|U|^2)(\operatorname{Re}(\overline U\eta))^2$.
The implication from negative energy to $v\cdot P<0$ must be checked
separately for such a model.  No such monotonicity is asserted for
every nonlinear correction in the preceding dilation cone.
\end{remark}

\begin{remark}[Positive-type Hartree interaction]
The full index bound also extends to the critical Hartree setting
with a positive interaction form $D$, under the Hartree form and
regularity assumptions used in this manuscript.  There the frozen
complex-linear operator still obeys $K_UU=0$ and
$K_U\overline U=2iv\cdot\nabla\overline U$, while
\[
 q_{\mathcal L_U}(\eta)=\langle\eta,K_U\eta\rangle_{\mathbb R}
 -2D\bigl(\operatorname{Re}(\overline U\eta),
          \operatorname{Re}(\overline U\eta)\bigr).
\]
Positivity of $D$ makes the correction nonpositive, and its value
on the amplitude direction is strictly negative when
$D(|U|^2,|U|^2)>0$.  The proof on
$\operatorname{span}_{\mathbb R}\{U,\overline U,i\overline U\}$
is unchanged.  A nonnegative subhomogeneous dispersion addition
preserves the implication from $E<0$ because its dilation identity
is $v\cdot P=2mE-D_b<0$.  Thus $n_-(\mathcal L_U)\ge3$
also holds in that class.  The conclusion uses positivity of the
interaction form, not merely pointwise nonnegativity of its kernel.
\end{remark}

\begin{corollary}[A continuation barrier from a ground state]
\label{cor:r10-continuation}
A real sharp mass-critical ground state has full action Hessian
negative index one.  A continuous Fredholm family of profiles starting
at such a state cannot reach a negative-energy profile while retaining
only its phase and translation kernel and keeping all other spectral
values uniformly separated from zero.  If the essential spectral gap
stays open, reaching a negative-energy profile requires a net gain of
at least two negative directions through additional zero eigenvalues,
counted with multiplicity.
\end{corollary}
\begin{proof}
At a real optimizer $Q$, the phase block is nonnegative.  On the real
mass tangent space $\{h:\langle Q,h\rangle=0\}$ the real block is
nonnegative as well: $E$ is nonnegative on the mass sphere
$M=M_c$ and vanishes at $Q$, so its constrained second variation is
nonnegative.  On $Q$ itself the real block has quadratic form
$-\alpha\|Q\|_{\alpha+2}^{\alpha+2}<0$.  Hence the full index is one.
The index is locally constant for continuous selfadjoint Fredholm
forms after quotienting a kernel of fixed dimension, as long as the
remaining spectrum is separated from zero.  Apply
Theorem~\ref{thm:r10-three-directions} at the proposed endpoint.
\end{proof}

\subsection{Additional negative directions from momentum response}

Assume here that $v\mapsto U_v$ is a $C^1$ family in $X$ of exact
profiles at fixed $\omega$, and that differentiating the profile
equation is justified in $X^*$.  Define
\[
 B=D_vP(U_v).
\]
The matrix is symmetric: if $d(v)=S_{\omega,v}(U_v)$, then
$d_v=-P(U_v)/2$ and $B=-2d_{vv}$.

\begin{theorem}[A response-refined index bound]
\label{thm:r10-response-index}
At a negative-energy member of this family,
\begin{equation}\label{eq:r10-response-index}
 n_-(\mathcal L_{U_v})\ge
 3+n_-\bigl(B|_{P(U_v)^\perp}\bigr).
\end{equation}
For an isotropic, rotationally equivariant family with
$P(U_v)=p(c)v/c$, $c=|v|>0$, this gives
\begin{equation}\label{eq:r10-isotropic-index}
 n_-(\mathcal L_{U_v})\ge N+2.
\end{equation}
Consequently such a wave is not a local energy minimum on the
joint mass and momentum level set, provided that level set is a
$C^2$ manifold near the wave.
\end{theorem}
\begin{proof}
Let $w_\zeta=D_vU_v[\zeta]$.  Profile differentiation gives
\[
 \mathcal L_{U_v}w_\zeta=-i\zeta\cdot\nabla U_v,
 \qquad
 \langle w_\eta,\mathcal L_{U_v}w_\zeta\rangle_{\mathbb R}
       =\tfrac12\eta^TB\zeta.
\]
The second formula uses
$DP(U)[w]=2\operatorname{Re}\langle w,-i\nabla U\rangle$.
The pairings of $\mathcal Lw_\zeta$ with $\overline U$ and
$i\overline U$ vanish because $\int U\partial_jU=0$.
Its pairing with $U$ is $\zeta\cdot P(U)$, which vanishes for
$\zeta\perp P(U)$.
Thus the image of a negative subspace of $B|_{P^\perp}$ is negative
and form-orthogonal to the three-dimensional space in
\eqref{eq:r10-three-plane}.  The map is injective on that subspace,
since its quadratic form is strictly negative.  The two negative
spaces have trivial intersection, proving \eqref{eq:r10-response-index}.

In the isotropic case, with $e=v/c$,
\[
 B=p'(c)e\otimes e+\frac{p(c)}c(I-e\otimes e).
\]
Negative energy gives $cp(c)=v\cdot P<0$, so $p(c)<0$ and
$B|_{P^\perp}=(p(c)/c)I$ has $N-1$ negative directions.
This proves \eqref{eq:r10-isotropic-index}.
The joint mass and momentum tangent space has codimension at most
$N+1$, so it intersects the $(N+2)$-dimensional negative subspace
nontrivially.  At a constrained critical point the corresponding
second variation of $E$ is the restriction of $\mathcal L$.
If the joint level set is a $C^2$ manifold, this negative tangent
direction is realized by a constraint-preserving curve and excludes
a local minimum.
\end{proof}

The last level-set assumption is not a hidden spectral hypothesis.
For a sufficiently regular nonzero profile its charge differentials
are independent: a real relation between them would give
$(\tau-i\zeta\cdot\nabla)U=0$, forcing its $L^2$ Fourier transform
onto a hyperplane (or forcing $U=0$ if $\zeta=0$).
The finite-dimensional implicit-function theorem then supplies the
local charge manifold in the energy space.

\subsection{An exact reduction of the charge Hessian}

Suppose now that $(\omega,v)\mapsto U_{\omega,v}$ is a $C^1$
profile family on an open parameter set, that
$d(\omega,v)=S_{\omega,v}(U_{\omega,v})$ is $C^2$, and that its
charges obey the critical scaling covariance
\begin{align*}
 M(U_{\lambda^m\omega,\lambda^{m-1}v})&=M(U_{\omega,v}),\\
 P(U_{\lambda^m\omega,\lambda^{m-1}v})&=\lambda P(U_{\omega,v})
\end{align*}
whenever the parameters are in the domain.  This follows, in
particular, if the family is closed under the usual critical dilation,
up to phase and translation.  Let $H=d''(\omega,v)$ and
$B=D_vP(U_{\omega,v})$.

\begin{theorem}[Bordered momentum formula]
\label{thm:r10-bordered}
If $\omega\ne0$, the symmetric matrices $H$ and
\begin{equation}\label{eq:r10-bordered}
 -\frac12
 \begin{pmatrix}
 (m-1)v\cdot P&P^T\\
 P&B
 \end{pmatrix}
\end{equation}
are congruent.  If $B$ is invertible, define
\begin{equation}\label{eq:r10-kappa}
 \kappa=P^TB^{-1}P-(m-1)v\cdot P.
\end{equation}
Then, writing $p(H)$ for the number of positive eigenvalues,
\begin{equation}\label{eq:r10-hessian-inertia}
 p(H)=n_-(B)+\mathbf1_{\{\kappa>0\}},
 \qquad \dim\ker H=\mathbf1_{\{\kappa=0\}}.
\end{equation}
The Schur complement of the $v$ block in $H$ is exactly
$\kappa/(2m^2\omega^2)$.
\end{theorem}
\begin{proof}
Stationarity gives $d_\omega=M/2$ and $d_v=-P/2$.
Put $D=(m\omega,(m-1)v)^T$.  Differentiating the two scaling laws
gives the exact column identity
\[
 HD=\begin{pmatrix}0\\-P/2\end{pmatrix},
 \qquad H_{vv}=-B/2.
\]
The invertible matrix
\[
 R=\begin{pmatrix}m\omega&0\\(m-1)v&I\end{pmatrix}
\]
therefore satisfies $R^THR=$ \eqref{eq:r10-bordered}.
Sylvester's law of inertia and the Schur complement of $B$ prove
\eqref{eq:r10-hessian-inertia}.
For the final assertion, the mixed block is
\[
 H_{v\omega}=-\frac{P-(m-1)Bv}{2m\omega},
\]
and the upper-left entry follows from $HD=(0,-P/2)^T$.
Substitution and cancellation give
$H_{\omega\omega}-H_{\omega v}H_{vv}^{-1}H_{v\omega}
=\kappa/(2m^2\omega^2)$.
\end{proof}

For an isotropic family $P=p(c)e$ with $p(c)p'(c)\ne0$,
\begin{equation}\label{eq:r10-elasticity}
 \kappa=\frac{p(c)^2}{p'(c)}-(m-1)cp(c).
\end{equation}
Thus charge-Hessian degeneracy occurs precisely at
$cp'(c)/p(c)=1/(m-1)$, under these nonvanishing assumptions.
This is a degeneracy criterion, not an assertion that every such
point changes dynamical stability.  The momentum and its derivative
at a fixed frequency suffice; a separate numerical differentiation
in frequency is unnecessary.

\subsection{The small-speed transition at dispersion order two}

Assume a $C^1$ branch at $\omega=1$ starts at a real standing
profile $Q$.  Assume $L_-=a(D)+1-|Q|^\alpha$ has kernel
$\mathbb RQ$ and a bounded inverse on its orthogonal complement.
Assume also that the real-part derivative of the branch at zero
vanishes after the chosen translation gauge.  One sufficient
condition for this last assertion is
$\ker L_+=\operatorname{span}\{\partial_jQ\}$ together with
the translation-orthogonal slice; neither this nondegeneracy nor
branch existence is asserted for all ground states.
Define
\[
 T_{ij}=\langle\partial_iQ,L_-^{-1}\partial_jQ\rangle.
\]
If $T$ has $k$ negative eigenvalues, then
\begin{equation}\label{eq:r10-standing-index}
 n_-(L_-)\ge k,\qquad
 n_-\bigl(\operatorname{diag}(L_+,L_-)\bigr)\ge2k+1.
\end{equation}
Indeed, on the inverse-response vectors $L_-^{-1}(\zeta\cdot\nabla Q)$
the $L_-$ form is exactly $\zeta^TT\zeta$.  Let $V$ be the resulting
$k$-dimensional negative subspace.  The form of $L_-$ is negative
on $V$ and vanishes on $Q$, with no cross term.  Since
$L_+=L_--\alpha|Q|^\alpha$, it is strictly negative on
$V\oplus\mathbb RQ$: for a nonzero $V$ component the first form is
already negative, and on a nonzero multiple of $Q$ the subtracted
potential term is strictly positive.  The two diagonal blocks then
give \eqref{eq:r10-standing-index}.
Finite strictly negative subspaces persist under form-continuous
perturbations, so this lower bound persists on sufficiently nearby
traveling profiles.  In particular an isotropic negative response
$T=I_0I_N$, $I_0<0$, requires at least $2N+1$ negative action
directions on the standing state and its nearby traveling branch.

Differentiating the imaginary equation gives
$D_{v_j}U|_{v=0}=-iL_-^{-1}\partial_jQ$, and hence
\begin{equation}\label{eq:r10-small-speed}
 P(U_{1,v})=2Tv+o(|v|),\qquad B(1,v)=2T+o(1).
\end{equation}
For a nonsingular $T$ and directions on which $v^TTv$ stays away
from zero relative to $|v|^2$, the preceding exact formulas yield
\begin{align}
 E(U_{1,v})&=\frac1m v^TTv+o(|v|^2),\label{eq:r10-energy-expansion}\\
 M(U_{1,v})&=M(Q)+\frac{m-2}{m}v^TTv+o(|v|^2),
 \label{eq:r10-mass-expansion}\\
 \kappa(1,v)&=2(2-m)v^TTv+o(|v|^2).
 \label{eq:r10-kappa-expansion}
\end{align}
For \eqref{eq:r10-mass-expansion}, integrate $d_v=-P/2$ to get
$d(1,v)=M(Q)/2-v^TTv/2+o(|v|^2)$ and use weighted homogeneity,
$d=d_\omega+(m-1)v\cdot d_v/m$ at $\omega=1$.
Only the stated $C^1$ branch and $C^2$ action regularity are used;
no fourth derivative of a noninteger power nonlinearity is assumed.

At a sharp ground state $T$ is positive definite by phase
coercivity.  For sufficiently small nonzero velocity the charge
Hessian is consequently negative definite if $m>2$, whereas it
has exactly one positive eigenvalue if $1<m<2$.
At $m=2$ this leading-order test vanishes.  General homogeneous
symbols of degree two need not be quadratic forms, so no exact
Galilean degeneracy is inferred for all such symbols.

These are charge and index statements.  Under the additional
hypotheses of a Hamiltonian index theorem---a closed Hamiltonian
linearization, finite negative index, a coercive complementary
subspace, exactly the phase and translation kernel, independent
charges, and an invertible charge Hessian---the standard index
formula reads
\begin{equation}\label{eq:r10-hamiltonian-index}
 k_r+2k_c+2k_i^-
 =n_-(\mathcal L)-p(H)
 =n_-(\mathcal L)-n_-(B)-\mathbf1_{\{\kappa>0\}}.
\end{equation}
Here the terms on the left count real unstable eigenvalues,
complex unstable pairs, and the negative index on nonzero
imaginary modes, respectively, with the conventions of the chosen
index theorem.  This last formula invokes the established
Hamiltonian index theory of Lin and Zeng~\cite{LZ}; it is not a replacement
for checking its operator hypotheses.  For example, if a
nondegenerate small-speed ground-state branch retains
$n_-(\mathcal L)=1$ and all those hypotheses hold, the $m>2$
case has index one and therefore one positive real eigenvalue.
No unconditional nonlinear stability or instability conclusion
is made for arbitrary branches.
